\subsection{Association Rules and Density-Ratio Views}

Association-rule analysis searches for simple descriptions of regions where the input distribution has relatively large probability.  The classical market-basket setting is the binary special case, but the same idea extends to conjunctions of coordinate conditions.

\begin{definition} \label{def:sl-conjunctive-region} \hlt{Conjunctive Region}\\
Let $X=(X_1,\ldots,X_p)$ take values in $\mathcal{S}=\prod_{j=1}^p\mathcal{S}_j$. For $\mathcal{J}\subseteq\{1,\ldots,p\}$ and $s=(s_j)_{j\in\mathcal{J}}$ with $s_j\subseteq\mathcal{S}_j$, the \hlt{conjunctive region} is $R(\mathcal{J},s)=\{x\in\mathcal{S}:x_j\in s_j\ \forall j\in\mathcal{J}\}$, with event $\{X\in R(\mathcal{J},s)\}=\bigcap\limits_{j\in\mathcal{J}}\{X_j\in s_j\}$.\\
Coordinates in $\{1,\ldots,p\}\setminus\mathcal{J}$ are unrestricted; under the extension $s_j=\mathcal{S}_j$, they add no constraint.\\
For $\mathcal{S}_j\subseteq\mathbb{R}$, $s_j$ is typically an interval; for categorical $\mathcal{S}_j$, $s_j$ is typically a subset of levels.
\end{definition}

\begin{definition} \label{def:sl-market-basket-itemset-support} \hlt{Market-Basket Item Sets}\\
Let $\mathcal{I}=\{1,\ldots,K\}$ and $z_i=(z_{ik})_{k\in\mathcal{I}}\in\{0,1\}^{\mathcal{I}}$ for $i=1,\ldots,n$, where $z_{ik}=1$ iff observation $i$ contains item $k$.
\begin{enumerate}[label=(\roman*)]
\setlength{\itemsep}{0pt}
\item An \hlt{item set} is any $\mathcal{K}\subseteq\mathcal{I}$.
\item The containment indicator is $Y_i(\mathcal{K})=\prod\limits_{k\in\mathcal{K}}z_{ik}$; observation $i$ contains $\mathcal{K}$ iff $Y_i(\mathcal{K})=1$.
\item The empirical support is $T(\mathcal{K})=n^{-1}\sum\limits_{i=1}^nY_i(\mathcal{K})$.
\item For finite categorical supports, take $\mathcal{I}=\{(j,a):1\leq j\leq p,\ a\in\mathcal{S}_j\}$, hence $K=\sum\limits_{j=1}^p\abs{\mathcal{S}_j}$.
\end{enumerate}
\end{definition}

\begin{definition} \label{def:sl-association-rule-confidence-lift} \hlt{Association Rule, Confidence, and Lift}\\
Let $A,B\subseteq\mathcal{I}$ satisfy $A\cap B=\emptyset$, $A\neq\emptyset$, $B\neq\emptyset$, $T(A)>0$, and $T(B)>0$.\\
The \hlt{association rule} with antecedent $A$ and consequent $B$ is $A\Rightarrow B$.
\begin{enumerate}[label=(\roman*)]
\setlength{\itemsep}{0pt}
\item The rule support is $S(A\Rightarrow B)=T(A\cup B)$.
\item The confidence is $C(A\Rightarrow B)=\frac{T(A\cup B)}{T(A)}$.
\item The lift is $L(A\Rightarrow B)=\frac{C(A\Rightarrow B)}{T(B)}=\frac{T(A\cup B)}{T(A)T(B)}$.
\end{enumerate}
Interpretation: $C(A\Rightarrow B)$ estimates prevalence of $B$ given $A$; $L(A\Rightarrow B)$ compares it with $T(B)$.
\end{definition}

\begin{example} \label{ex:sl-market-basket-support-confidence-lift} \hlt{A Market-Basket Rule}\\
Let $A=\{\mathrm{peanut\ butter},\mathrm{jelly}\}$ and $B=\{\mathrm{bread}\}$.
\begin{enumerate}[label=(\roman*)]
\setlength{\itemsep}{0pt}
\item Support $T(A\cup B)=0.03$ means all three items appear in about $3\%$ of baskets.
\item Confidence $C(A\Rightarrow B)=0.82$ means about $82\%$ of baskets containing $A$ also contain bread.
\item If $T(B)=0.43$, then $L(A\Rightarrow B)=0.82/0.43\approx1.91$.
\end{enumerate}
The rule predicts bread more often than its marginal prevalence.
\end{example}

\begin{breakablealgorithm}
\caption{Apriori Association-Rule Search}
\label{alg:sl-apriori-association-rules}
\begin{algorithmic}
\Require Binary item data $z_1,\ldots,z_n$, support threshold $t$, confidence threshold $c$
\Ensure Rules with support above $t$ and confidence above $c$
\State Compute supports for all one-item sets and retain $\mathcal{F}_1=\{\mathcal{K}:|\mathcal{K}|=1,\ T(\mathcal{K})>t\}$
\State Set $m\gets2$
\While{$\mathcal{F}_{m-1}\neq\emptyset$}
\State Form candidate item sets of size $m$ whose every size-$(m-1)$ subset lies in $\mathcal{F}_{m-1}$
\State Scan the data to compute $T(\mathcal{K})$ for each candidate
\State Retain $\mathcal{F}_m=\{\mathcal{K}:|\mathcal{K}|=m,\ T(\mathcal{K})>t\}$
\State Set $m\gets m+1$
\EndWhile
\State For each retained item set $\mathcal{K}$, test nonempty proper antecedents $A\subset\mathcal{K}$
\State Report $A\Rightarrow \mathcal{K}\setminus A$ when $C(A\Rightarrow \mathcal{K}\setminus A)>c$
\end{algorithmic}
\end{breakablealgorithm}

\begin{remark} \label{rem:sl-apriori-monotonicity-and-limits} \hlt{Apriori Monotonicity and Limits}\\
Support is antitone: $\mathcal{L}\subseteq\mathcal{K}\Longrightarrow T(\mathcal{L})\geq T(\mathcal{K})$.
\begin{enumerate}[label=(\roman*)]
\setlength{\itemsep}{0pt}
\item If $T(\mathcal{L})<t$, then every $\mathcal{K}\supseteq\mathcal{L}$ also has $T(\mathcal{K})<t$ and may be pruned.
\item A support threshold $t$ can discard rules $A\Rightarrow B$ with large $C(A\Rightarrow B)$ or $L(A\Rightarrow B)$ but small $T(A\cup B)$.
\end{enumerate}
\end{remark}

\begin{figure}[H]
\centering
\begin{tikzpicture}[
  >=Stealth,
  every node/.style={font=\scriptsize},
  item/.style={draw=gray!65, rounded corners=2pt, minimum width=0.72cm, minimum height=0.42cm, inner sep=1pt, fill=white},
  keep/.style={item, draw=green!55!black, fill=green!9},
  prune/.style={item, draw=red!70!black, fill=red!8},
  edge/.style={gray!45, line width=0.35pt},
  badedge/.style={red!70!black, dashed, line width=0.65pt},
  labelbox/.style={fill=white, inner sep=1.2pt},
  level/.style={text=gray!70!black}
]
  \path[use as bounding box] (-0.55,0.45) rectangle (5.45,3.35);

  \node[level] at (-0.30,3.0) {$m=1$};
  \node[level] at (-0.30,1.95) {$m=2$};
  \node[level] at (-0.30,0.90) {$m=3$};

  \node[keep] (A) at (0.75,3.0) {$A$};
  \node[keep] (B) at (1.85,3.0) {$B$};
  \node[keep] (C) at (2.95,3.0) {$C$};
  \node[keep] (D) at (4.05,3.0) {$D$};

  \node[keep] (AC) at (0.95,1.95) {$AC$};
  \node[prune] (AB) at (1.85,1.95) {$AB$};
  \node[keep] (BC) at (2.95,1.95) {$BC$};
  \node[keep] (CD) at (4.05,1.95) {$CD$};

  \node[prune] (ABC) at (1.85,0.90) {$ABC$};
  \node[prune] (ABD) at (2.75,0.90) {$ABD$};
  \node[keep] (BCD) at (3.85,0.90) {$BCD$};

  \foreach \u/\v in {A/AC,A/AB,B/AB,B/BC,C/AC,C/BC,C/CD,D/CD,BC/BCD,CD/BCD}
    \draw[edge] (\u) -- (\v);
  \draw[badedge] (AB) -- (ABC);
  \draw[badedge] (AB) -- (ABD);

\end{tikzpicture}

\caption{Apriori keeps \textcolor{green!55!black}{frequent item sets}; if \textcolor{red!70!black}{$T(AB)\le t$}, all supersets of $AB$ are pruned.}
\label{fig:sl-apriori-itemset-monotonicity}
\end{figure}

\begin{method} \label{meth:sl-generalized-association-rules} \hlt{Generalized Association Rules}\\
For $\mathcal{J}\subseteq\{1,\ldots,p\}$ and $s_j\subseteq\mathcal{S}_j$, set $R(\mathcal{J},s)=\bigcap\limits_{j\in\mathcal{J}}\{X_j\in s_j\}$ and $\widehat{P}\{R(\mathcal{J},s)\}=n^{-1}\sum\limits_{i=1}^n\mathbf{1}\{x_{ij}\in s_j\ \forall j\in\mathcal{J}\}$.
\begin{enumerate}[label=(\roman*)]
\setlength{\itemsep}{0pt}
\item Allow $s_j$ to be intervals for quantitative variables and subsets of levels for categorical variables.
\item Search over candidate regions with a heuristic rule learner.
\item Retain rules whose fitted score and observed support meet the chosen thresholds.
\end{enumerate}
\end{method}

\begin{method} \label{meth:sl-unsupervised-as-supervised-density-ratio} \hlt{Unsupervised Learning as Supervised Density-Ratio Estimation}\\
Let $g$ be the data density, $g_0$ a reference density, and $Y=1$ label data while $Y=0$ labels reference observations.
\begin{enumerate}[label=(\roman*)]
\setlength{\itemsep}{0pt}
\item Draw a reference sample from $g_0$ and pool it with the data, weighting the two classes equally.
\item Fit a probabilistic classifier $\widehat{\mu}(x)\approx\Pr(Y=1\mid X=x)$.
\item Use $\mu(x)=g(x)/(g(x)+g_0(x))$ to estimate the density ratio by $\widehat{g}(x)/g_0(x)=\widehat{\mu}(x)/(1-\widehat{\mu}(x))$.
\end{enumerate}
Uniform, Gaussian, and product-marginal $g_0$ target high density, nonnormality, and dependence, respectively.
\end{method}

\begin{remark} \label{rem:sl-rule-learners-for-density-ratios} \hlt{Rule Learners for Density Ratios}\\
After estimating $\widehat{g}/g_0$, any learner that returns regions $R=\bigcap\limits_{j\in\mathcal{J}}\{X_j\in s_j\}$ can search for generalized item sets. Tree leaves and box-peeling regions have this conjunctive form.\\
The unsupervised target is large $\widehat{g}(x)/g_0(x)$ on $R$ together with adequate empirical support, not predictive risk.
\end{remark}
